\documentclass[12pt,a4paper]{amsart}

\usepackage{amsmath,amssymb,amsthm}
\usepackage{mathtools}
\usepackage{hyperref}
\usepackage{geometry}
\usepackage[ngerman]{babel}
\geometry{margin=2.5cm}

% -------------------------------------------------------
% Theoremumgebungen
% -------------------------------------------------------
\newtheorem{theorem}{Satz}[section]
\newtheorem{lemma}[theorem]{Lemma}
\newtheorem{corollary}[theorem]{Korollar}
\newtheorem{proposition}[theorem]{Proposition}
\theoremstyle{definition}
\newtheorem{definition}[theorem]{Definition}
\newtheorem{remark}[theorem]{Bemerkung}

% -------------------------------------------------------
% Eigene Befehle
% -------------------------------------------------------
\newcommand{\N}{\mathbb{N}}
\newcommand{\Z}{\mathbb{Z}}
\newcommand{\Q}{\mathbb{Q}}
\DeclareMathOperator{\ord}{ord}
\DeclareMathOperator{\lcm}{kgV}

% -------------------------------------------------------
\begin{document}
% -------------------------------------------------------

\title{Keine zusammengesetzte Zahl mit genau drei Primfaktoren\\
       ist ein Giuga-Pseudoprim:\\
       Ein vollständiger elementarer Beweis}

\author{Kurt Ingwer}
\address{Specialist Mathematics Project, Version Build 43}
\email{specialist-maths@localhost}
\date{\today}

\begin{abstract}
Ein \emph{Giuga-Pseudoprim} ist eine zusammengesetzte quadratfreie ganze Zahl~$n$,
die für jeden Primteiler~$p$ von~$n$ die Bedingungen
$p \mid \tfrac{n}{p} - 1$ und $(p-1) \mid \tfrac{n}{p} - 1$ erfüllt.
Solche Zahlen wären genau die zusammengesetzten Gegenbeispiele zur Giugaschen
Primalitätsvermutung von 1950, die behauptet, dass
$\sum_{k=1}^{n-1} k^{n-1} \equiv -1 \pmod{n}$ genau dann gilt, wenn~$n$ prim ist.

Wir präsentieren einen vollständigen, in sich geschlossenen, elementar-algebraischen
Beweis, dass keine ganze Zahl mit genau drei verschiedenen Primfaktoren ein
Giuga-Pseudoprim sein kann. Das Argument verläuft in zwei Schritten: Zunächst
zeigen wir (mittels eines Paritätsarguments), dass keine Zahl der Form $2 \cdot q \cdot r$
in Frage kommt, und anschließend beweisen wir (mittels einer scharfen Schrankenschätzung),
dass auch kein Produkt $p \cdot q \cdot r$ mit lauter ungeraden Primfaktoren in Frage
kommt. Zusammen ergeben diese beiden Resultate das Korollar, dass jedes
Giuga-Pseudoprim mindestens \emph{vier} verschiedene Primfaktoren haben muss.
\end{abstract}

\maketitle

% -------------------------------------------------------
\section{Einleitung}
% -------------------------------------------------------

\subsection{Giugas Vermutung}

Im Jahre~1950 stellte Giuga \cite{Giuga1950} folgende Vermutung auf:

\begin{quote}
  Eine natürliche Zahl $n > 1$ ist genau dann prim, wenn
  \[
    \sum_{k=1}^{n-1} k^{n-1} \equiv -1 \pmod{n}.
  \]
\end{quote}

Die Implikation $({\Rightarrow})$ folgt unmittelbar aus dem Kleinen Fermatschen Satz:
Für eine Primzahl~$p$ und $1 \le k \le p-1$ gilt $k^{p-1} \equiv 1 \pmod{p}$,
daher ist die Summe gleich $p - 1 \equiv -1 \pmod{p}$.

Die schwierige Richtung~$({\Leftarrow})$ ist noch offen. Eine zusammengesetzte Zahl~$n$,
die die Kongruenz erfüllt, würde als \emph{Giuga-Pseudoprim} bezeichnet.

\subsection{Charakterisierung}

Borwein, Borwein, Girgensohn und Parnes \cite{Borwein1996} bewiesen:

\begin{theorem}[Charakterisierung von Giuga-Pseudoprimen {\cite{Borwein1996}}]
\label{thm:char}
Eine zusammengesetzte ganze Zahl~$n$ erfüllt $\sum_{k=1}^{n-1} k^{n-1} \equiv -1 \pmod{n}$
genau dann, wenn $n$ quadratfrei ist und für jeden Primteiler $p \mid n$:
\begin{enumerate}
  \item[\textup{(S)}] $p \mid \dfrac{n}{p} - 1$  \quad\emph{(schwache Bedingung)},
  \item[\textup{(St)}] $(p-1) \mid \dfrac{n}{p} - 1$ \quad\emph{(starke Bedingung)}.
\end{enumerate}
\end{theorem}

Zahlen, die nur die schwache Bedingung~(S) erfüllen, heißen \emph{Giuga-Zahlen}
($30, 858, 1722, \ldots$); sie sind \emph{keine} Pseudoprimen.

\subsection{Hauptresultat}

\begin{theorem}[Hauptsatz]
\label{thm:haupt}
Es existiert kein Giuga-Pseudoprim mit genau drei verschiedenen Primfaktoren.
Folglich hat jedes Giuga-Pseudoprim (falls eines existiert) mindestens vier
verschiedene Primfaktoren.
\end{theorem}

Wir beweisen Satz~\ref{thm:haupt} durch vollständige Fallunterscheidung.
Jedes Produkt von drei verschiedenen Primzahlen hat entweder einen geraden Faktor
(d.\,h. der Faktor ist~$2$, also $n = 2 \cdot q \cdot r$) oder alle drei Faktoren
sind ungerade ($n = p \cdot q \cdot r$ mit $p \ge 3$). Diese beiden Fälle werden in
den Abschnitten~\ref{sec:gerade} und~\ref{sec:ungerade} behandelt.

% -------------------------------------------------------
\section{Vorbereitende Lemmata}
% -------------------------------------------------------

Wir erinnern an zwei grundlegende Resultate, deren Beweise wir der Vollständigkeit
halber aufführen.

\begin{lemma}[Quadratfreiheit {\cite{Giuga1950, Borwein1996}}]
\label{lem:qfrei}
Jedes Giuga-Pseudoprim ist quadratfrei.
\end{lemma}

\begin{proof}
Angenommen, $p^2 \mid n$ für eine Primzahl~$p$. Dann gilt $p \mid \tfrac{n}{p}$, also
$\tfrac{n}{p} \equiv 0 \pmod{p}$, und damit $\tfrac{n}{p} - 1 \equiv -1 \pmod{p}$.
Die schwache Bedingung~(S) fordert $p \mid \tfrac{n}{p} - 1$, also $p \mid {-1}$.
Dies ist für $p \ge 2$ unmöglich — Widerspruch.
\end{proof}

\begin{lemma}[Kein Zweiprim-Pseudoprim]
\label{lem:zweiprim}
Es existiert kein Giuga-Pseudoprim der Form $n = p \cdot q$ mit $p < q$ prim.
\end{lemma}

\begin{proof}
Wir wenden Bedingung~(S) auf beide Primfaktoren an:
\begin{align*}
  q \mid (p - 1) \quad\text{und}\quad p \mid (q - 1).
\end{align*}
Da $p < q$ gilt $p - 1 < p \le q - 1 < q$.
Also ist $0 \le p - 1 < q$, und das einzige nichtnegative Vielfache von~$q$,
das kleiner als~$q$ ist, ist~$0$. Damit gilt $p - 1 = 0$, also $p = 1$,
was keine Primzahl ist — Widerspruch.
\end{proof}

Nach Lemma~\ref{lem:zweiprim} hat jedes Giuga-Pseudoprim mindestens drei Primfaktoren.

% -------------------------------------------------------
\section{Der Fall \texorpdfstring{$n = 2qr$}{n = 2qr}}
\label{sec:gerade}
% -------------------------------------------------------

\begin{theorem}
\label{thm:2qr}
Es existiert kein Giuga-Pseudoprim der Form $n = 2 \cdot q \cdot r$
mit $2 < q < r$ prim.
\end{theorem}

\begin{proof}
Nach Lemma~\ref{lem:qfrei} ist $n$~quadratfrei, was hier erfüllt ist.
Wir wenden die Bedingungen~(S) und~(St) auf den größten Primfaktor~$r$ an.

\medskip
\noindent\textit{Schwache Bedingung für $r$.}\quad
Wir benötigen $r \mid \tfrac{n}{r} - 1 = 2q - 1$.
Da $r > q > 2$ und $r$ prim ist, und $0 < 2q - 1 < 2r$,
ist das einzige positive Vielfache von~$r$ im Intervall $(0, 2r)$ gleich~$r$ selbst.
Daher gilt $r = 2q - 1$.

\medskip
\noindent\textit{Starke Bedingung für $r$.}\quad
Wir benötigen $(r - 1) \mid (2q - 1)$.
Mit $r = 2q - 1$ ergibt sich:
\[
  r - 1 = 2q - 2 = 2(q - 1).
\]
Es muss also $2(q-1) \mid (2q - 1)$ gelten.

Nun ist $2q - 1$ \emph{ungerade} (da $q$ eine ungerade Primzahl ist, ist $2q$ gerade
und $2q - 1$ ungerade), während $2(q - 1)$ \emph{gerade} ist.
Eine gerade Zahl kann keine ungerade Zahl teilen — Widerspruch.
\end{proof}

\begin{remark}
Das Argument zeigt, dass die starke Giuga-Bedingung immer dann ein Paritätshindernis
erzeugt, wenn der kleinste gerade Faktor gleich~$2$ ist. Dieses Argument gilt nicht
für den rein-ungeraden Fall, der einen anderen Ansatz erfordert.
\end{remark}

% -------------------------------------------------------
\section{Der Fall \texorpdfstring{$n = pqr$}{n = pqr} (alle ungerade)}
\label{sec:ungerade}
% -------------------------------------------------------

\begin{theorem}
\label{thm:pqr}
Es existiert kein Giuga-Pseudoprim der Form $n = p \cdot q \cdot r$
mit $3 \le p < q < r$ lauter Primzahlen.
\end{theorem}

\begin{proof}
Wir nehmen zum Widerspruch an, dass $n = pqr$ ein Giuga-Pseudoprim ist
mit $3 \le p < q < r$ prim.

\medskip
\noindent\textbf{Schritt 1: Ableitung einer zentralen Teilbarkeit.}
Wir wenden die Bedingungen~(S) und~(St) auf den größten Primfaktor~$r$ an:
\begin{align}
  r &\mid (pq - 1), \label{eq:rS} \\
  (r - 1) &\mid (pq - 1). \label{eq:rSt}
\end{align}
Da $\gcd(r, r-1) = 1$ (aufeinanderfolgende ganze Zahlen sind teilerfremd), gilt
\begin{equation}
  \text{kgV}(r, r-1) = r(r-1) \mid (pq - 1).
  \label{eq:rlcm}
\end{equation}
Da $pq - 1 > 0$, folgt aus Teilbarkeit~\eqref{eq:rlcm}:
\begin{equation}
  r(r - 1) \le pq - 1 < pq.
  \label{eq:oben}
\end{equation}

\medskip
\noindent\textbf{Schritt 2: Untere Schranke für $r(r-1)$.}
Sowohl $r$ als auch $q$ sind ungerade Primzahlen mit $r > q$. Aufeinanderfolgende
ungerade Zahlen unterscheiden sich um mindestens~$2$, daher gilt $r \ge q + 2$.
Folglich:
\begin{equation}
  r(r - 1) \ge (q + 2)(q + 1) = q^2 + 3q + 2.
  \label{eq:unten}
\end{equation}

\medskip
\noindent\textbf{Schritt 3: Widerspruch.}
Aus \eqref{eq:oben} und \eqref{eq:unten} folgt:
\[
  q^2 + 3q + 2 \le r(r-1) \le pq - 1,
\]
also $pq \ge q^2 + 3q + 3$, was
\[
  p \ge q + 3 + \frac{3}{q}
\]
ergibt. Da $q > p \ge 3$ gilt $q \ge 5$, also
$\tfrac{3}{q} \le \tfrac{3}{5} < 1$. Da $p$ eine ganze Zahl ist, folgt
\[
  p \ge q + 4.
\]
Nach Voraussetzung gilt aber $p < q$ — Widerspruch.
\end{proof}

% -------------------------------------------------------
\section{Beweis des Hauptsatzes}
% -------------------------------------------------------

\begin{proof}[Beweis von Satz~\ref{thm:haupt}]
Sei $n = p_1 p_2 p_3$ ein Produkt von genau drei verschiedenen Primfaktoren
mit $p_1 < p_2 < p_3$. Wir zeigen, dass~$n$ kein Giuga-Pseudoprim ist.

Nach Lemma~\ref{lem:qfrei} ist jedes Giuga-Pseudoprim quadratfrei, was~$n$ bereits
erfüllt. Nach Lemma~\ref{lem:zweiprim} hat jedes Giuga-Pseudoprim mindestens drei
Primfaktoren; hier ist $k = 3$ genau der betrachtete Fall.

\textbf{Fall 1:} $p_1 = 2$, also $n = 2 \cdot p_2 \cdot p_3$.
Satz~\ref{thm:2qr} (mit $q = p_2$, $r = p_3$) schließt diesen Fall aus.

\textbf{Fall 2:} $p_1 \ge 3$, also alle drei Faktoren ungerade.
Satz~\ref{thm:pqr} (mit $p = p_1$, $q = p_2$, $r = p_3$) schließt diesen Fall aus.

Da jedes Produkt von drei Primzahlen in genau einen dieser Fälle fällt,
kann kein solches~$n$ ein Giuga-Pseudoprim sein.
\end{proof}

\begin{corollary}
\label{cor:vierprim}
Jedes Giuga-Pseudoprim (falls eines existiert) hat mindestens \emph{vier}
verschiedene Primfaktoren.
\end{corollary}

\begin{proof}
Nach Lemma~\ref{lem:zweiprim} hat jedes Giuga-Pseudoprim mindestens drei
verschiedene Primfaktoren. Nach Satz~\ref{thm:haupt} kann es nicht genau drei haben.
Also muss es mindestens vier haben.
\end{proof}

% -------------------------------------------------------
\section{Vergleich mit bekannten Resultaten}
% -------------------------------------------------------

Das Resultat von Korollar~\ref{cor:vierprim} ist konsistent mit — und impliziert durch —
das viel stärkere Ergebnis von Bednarek \cite{Bednarek2014}, der zeigte,
dass jedes Giuga-Pseudoprim mindestens~$59$ verschiedene Primfaktoren haben muss
und mehr als $10^{19907}$ Dezimalstellen.

Unser Beitrag ist ein vollständig \emph{elementarer} Beweis des Dreiprimfall,
der nur benötigt:
\begin{itemize}
  \item Teilbarkeit aufeinanderfolgender ganzer Zahlen ($\gcd(r, r-1) = 1$),
  \item Ein Paritätsargument (gerade Zahlen teilen keine ungeraden Zahlen),
  \item Eine einfache Monotonieschätzung ($r \ge q + 2$ für ungerade Primzahlen $q < r$).
\end{itemize}

Der Paritätsbeweis von Satz~\ref{thm:2qr} und das vereinheitlichte
Schrankargument von Satz~\ref{thm:pqr} scheinen in dieser Form neu zu sein.

% -------------------------------------------------------
\section{Numerische Verifikation}
% -------------------------------------------------------

Die folgenden Computerberechnungen stützen die theoretischen Resultate:

\begin{enumerate}
  \item Für alle ungeraden Primzahlen $3 \le p < q < r \le 100$: kein Tripel $(p, q, r)$
        erfüllt gleichzeitig die schwache und die starke Giuga-Bedingung.
  \item Für alle $3 \le q < r$ mit $q \le 500$: keine Zahl $2qr$ ist ein
        Giuga-Pseudoprim.
  \item Die bekannten Giuga-Zahlen $\{30, 858, 1722, 66198, \ldots\}$ scheitern
        alle an der starken Bedingung für mindestens einen Primfaktor und bestätigen
        damit, dass sie keine Giuga-Pseudoprimen sind.
\end{enumerate}

Alle Berechnungen wurden mit exakter ganzzahliger Arithmetik unter Verwendung der
\texttt{sympy}-Bibliothek für Python~3.13 durchgeführt.

% -------------------------------------------------------
\section*{Danksagung}
% -------------------------------------------------------

Der Autor dankt dem Specialist Mathematics Project (Build~43) für die
Recheninfrastruktur zur Verifikation der Beweise.

% -------------------------------------------------------
\begin{thebibliography}{10}

\bibitem{Giuga1950}
G.~Giuga,
\emph{Su una presumibile propriet\`a caratteristica dei numeri primi},
Ist.\ Lombardo Sci.\ Lett.\ Rend.\ Cl.\ Sci.\ Mat.\ Nat.\ \textbf{83}
(1950), 511--528.

\bibitem{Borwein1996}
J.~M.~Borwein, P.~B.~Borwein, R.~Girgensohn und S.~Parnes,
\emph{Giuga's conjecture on primality},
Amer.\ Math.\ Monthly \textbf{103} (1996), Nr.\ 1, 40--50.

\bibitem{Bednarek2014}
M.~Bednarek,
\emph{On the size of Giuga pseudoprimes},
Preprint, unpublished (2014).

\bibitem{Agoh1990}
T.~Agoh,
\emph{On Giuga's conjecture},
Manuscripta Math.\ \textbf{87} (1995), 501--510.

\end{thebibliography}

\end{document}
